\section{逻辑公理系统}

\subsection{命题逻辑公理系统}

命题逻辑的公理系统包含如下要素：

\begin{itemize}
	\item 符号集合：
	\begin{itemize}
		\item 命题变元 $p_1, p_2$ 等；
		\item 联结词符号 $\lnot, \to$ 等；
		\item 括号 $(, )$。
	\end{itemize}
	\item 公式的形成规则：
	\begin{itemize}
		\item 若 $Q$ 是公式，那么 $(Q)$ 是公式；
		\item 若 $Q$ 是公式，那么 $\lnot Q$ 是公式；
		\item 若 $Q, R$ 是公式，那么 $Q \to R$ 是公式。
	\end{itemize}
	\item 公理：设 $P, Q, R$ 为任意公式，那么
	\begin{itemize}
		\item 公理模式 $\mathcal{A}_1$（肯定后件律）：$Q \to (R \to Q)$；
		\item 公理模式 $\mathcal{A}_2$（蕴含词分配律）：$(P \to (Q \to R)) \to ((P \to Q) \to (P \to R))$；
		\item 公理模式 $\mathcal{A}_3$（换位律）：$(\lnot Q \to \lnot R) \to (R \to Q)$。
	\end{itemize}
	\item 推理规则：份力规则，Modus Ponents 规则：
	\begin{itemize}
		\item 若 $Q$ 和 $Q \to R$ 成立，那么 $R$ 成立。这里 $Q$ 和 $Q \to R$ 称为前提，$R$ 称为结论；
	\end{itemize}
	\item 定理集：由公理和推理规则唯一确定出来的集合。
\end{itemize}

在命题公理系统中仅使用了 $\lnot$ 和 $\to$。其它联结词可以从这两个联结词定义出来。